\section{Repaso práctico y plantillas de casos}
\subsection{Ejemplo de corrección de código}
Graham describió un caso típico: después de que el Agent encuentra un bug, en la etapa de Plan se escribe la lista de archivos a modificar, en la etapa de Act se usa \texttt{python apply\_patch} para generar el diff, y en la etapa de Verify se ejecuta \texttt{pytest}. Todo el proceso queda registrado en el log, y si la prueba falla se revierte automáticamente.

\begin{knowledgebox}{Ejemplo de flujo}
Plan: extraer el Issue, listar archivos y funciones; Act: ejecutar el Agent para escribir el patch y guardar el diff; Verify: ejecutar las pruebas, y si fallan, enviar el log al issue.
\end{knowledgebox}

\subsection{Registro de fallas y correcciones}
El equipo registra cada muestra de falla en una «matriz de fallas»: incluye la causa de la falla (prueba, permisos, sintaxis), la salida del Agent y el contenido de la intervención humana. Esta matriz es material importante para entrenar los prompts.

\begin{warningbox}{No borren el log de fallas}
El log de fallas contiene información semántica clave, que se usa para generar datos de entrenamiento de agent chains; borrarlo impediría reproducir la falla en el futuro.
\end{warningbox}

\subsection{Resumen de la sección}
\begin{itemize}
    \item Mediante un flujo con plantilla, un success individual se puede definir como la conclusión de las tres etapas «Plan + Act + Verify».
    \item Registrar la matriz de fallas ofrece una base para el entrenamiento posterior y el ajuste de prompts.
\end{itemize}
